\[ \section{Introduction} The accurate prediction of nonlinear dynamical systems over extended temporal horizons remains one of the fundamental challenges in modern scientific computing. From computational fluid dynamics and plasma physics to climate modeling, structural mechanics, reaction--diffusion systems, and digital twins of engineering infrastructures, numerous scientific and engineering disciplines rely on computational models capable of predicting the evolution of complex physical processes. These systems are typically governed by nonlinear ordinary or partial differential equations whose analytical solutions are unavailable, while high-fidelity numerical simulations often require substantial computational resources that limit their applicability in real-time decision making, optimization, and control. Recent advances in Scientific Machine Learning (SciML) have demonstrated that deep neural networks can significantly accelerate physical simulations by learning surrogate representations of complex dynamical processes directly from numerical or experimental data. Neural operators, graph neural networks, physics-informed neural networks, differentiable simulators, and operator-learning architectures have substantially expanded the range of computational techniques available for modeling high-dimensional physical systems. These approaches have shown remarkable success in approximating solution operators for nonlinear partial differential equations while reducing computational costs by several orders of magnitude compared with conventional numerical solvers. Despite these advances, reliable long-horizon prediction remains an unresolved problem. Most existing learning-based models exhibit satisfactory performance during short-term forecasting but gradually accumulate approximation errors during autoregressive inference. Even small prediction inaccuracies introduced at early time steps may propagate nonlinearly through the temporal evolution of the system, ultimately producing physically inconsistent trajectories, unstable numerical behavior, and divergence from the underlying governing dynamics. This phenomenon becomes particularly severe in chaotic systems, strongly coupled multiphysics simulations, turbulent flows, and nonlinear reaction networks, where local perturbations may grow exponentially over time. One of the principal limitations of purely data-driven learning is that optimization objectives are primarily constructed to minimize prediction error on the available training data rather than to preserve the physical structure of the underlying dynamical process. Consequently, conventional neural architectures frequently violate conservation laws, symmetry properties, geometric invariants, and other physically meaningful constraints that naturally arise from the governing equations. Although these violations may remain negligible during one-step prediction, they often become the dominant source of instability during recursive long-term simulation. The emergence of Physics-Informed Neural Networks (PINNs) represented an important milestone toward integrating physical knowledge into machine learning. By incorporating residuals of governing differential equations into the optimization objective, PINNs demonstrated that data-driven learning can be regularized using first-principles physical models. Nevertheless, practical applications have revealed several limitations, including optimization difficulties for stiff systems, sensitivity to hyperparameter selection, challenges associated with high-dimensional domains, and reduced scalability for large spatiotemporal simulations. Moreover, PINNs primarily operate on continuous function approximation and are not inherently designed to exploit the relational structure naturally present in discretized physical systems. Graph Neural Networks (GNNs) provide an alternative computational paradigm in which physical systems are represented as attributed graphs whose nodes correspond to spatial entities and whose edges describe local physical interactions. Such graph representations naturally accommodate irregular computational meshes, particle methods, finite-volume discretizations, molecular systems, and heterogeneous physical domains. Message passing mechanisms enable efficient information exchange between neighboring elements while preserving local topological relationships, making graph neural architectures particularly attractive for scientific simulations involving complex geometries. However, existing graph-based surrogate models remain predominantly data-centric. In most architectures, physical knowledge enters the learning process only indirectly through training data rather than through explicit computational mechanisms that continuously enforce physically meaningful behavior during prediction. As a consequence, graph-based models frequently experience trajectory drift, loss of conservation properties, degradation of numerical stability, and increasing prediction uncertainty as forecasting horizons become longer. These observations motivate the development of a new generation of scientific learning architectures that integrate graph intelligence and physical reasoning at the architectural level rather than treating physics as an auxiliary regularization term. Such integration requires computational mechanisms capable of jointly modeling graph topology, temporal dynamics, governing physical principles, uncertainty propagation, and numerical stability within a unified differentiable framework. \] 